\corrigesChapitre{16}{Résistance et loi d’Ohm}
\par\noindent\begin{minipage}{\linewidth}
\begin{corrige}{1}{L’unité de R}
\begin{questions}
\item L’ohm.
\item $\Omega$.
\end{questions}
\end{corrige}
\end{minipage}
\par\noindent\begin{minipage}{\linewidth}
\begin{corrige}{2}{La relation}
\begin{questions}
\item $U=RI$.
\item U en volts, R en ohms et I en ampères.
\end{questions}
\end{corrige}
\end{minipage}
\par\noindent\begin{minipage}{\linewidth}
\begin{corrige}{3}{Convertir}
\begin{questions}
\item $0{,}040\unite A$.
\item $2\,200\unite\Omega$.
\end{questions}
\end{corrige}
\end{minipage}
\par\noindent\begin{minipage}{\linewidth}
\begin{corrige}{4}{Une tension directe}
\begin{questions}
\item $U=100\times0{,}02=2\unite V$.
\item Oui avec les unités de la loi.
\end{questions}
\end{corrige}
\end{minipage}
\par\noindent\begin{minipage}{\linewidth}
\begin{corrige}{5}{Une droite}
\begin{questions}
\item Une droite.
\item Oui : U est proportionnelle à I dans le domaine étudié.
\end{questions}
\pcFigureEx{figures/cor-ohm-modele.pdf}{La proportionnalité se reconnaît à la droite passant par l’origine, dans les conditions du modèle.}
\end{corrige}
\end{minipage}
\par\noindent\begin{minipage}{\linewidth}
\begin{corrige}{6}{Retrouver le courant}
\begin{questions}
\item $I=6/200=0{,}030\unite A=30\unite{mA}$.
\item $200\times0{,}030=6\unite V$.
\end{questions}
\par\smallskip\noindent\textbf{Pour comprendre.}\enspace De $U=RI$, isoler $I=U/R=6/200=0{,}030\unite A=30\unite{mA}$. Vérifier : $200\times0{,}030=6\unite V$.\par
\end{corrige}
\end{minipage}
\par\noindent\begin{minipage}{\linewidth}
\begin{corrige}{7}{Retrouver R}
\begin{questions}
\item $I=0{,}020\unite A$ ; $R=4{,}0/0{,}020=200\unite\Omega$.
\item Elle provient de l’emploi de 20 au lieu de 0,020 pour une intensité en ampères.
\end{questions}
\par\smallskip\noindent\textbf{Pour comprendre.}\enspace Convertir $20\unite{mA}=0{,}020\unite A$, puis $R=U/I=4{,}0/0{,}020=200\unite\Omega$. Le calcul 4/20 oublie le changement d’unité et donne un résultat 1 000 fois trop petit.\par
\end{corrige}
\end{minipage}
\par\noindent\begin{minipage}{\linewidth}
\begin{corrige}{8}{Une table cohérente}
\begin{questions}
\item Les quotients valent tous $150\unite\Omega$.
\item $U=150\times0{,}04=6\unite V$, si le domaine et les caractéristiques du composant le permettent.
\end{questions}
\par\smallskip\noindent\textbf{Pour comprendre.}\enspace Les quotients $1{,}5/0{,}01$, $3/0{,}02$ et $4{,}5/0{,}03$ valent tous 150 ohms. Prévoir ensuite $U=150\times0{,}04=6\unite V$, sans confondre lecture d’un point et coefficient de la droite.\par
\pcFigureEx{figures/cor-ohm-150.pdf}{Les points donnent 150 ohms ; le point prévu pour 0,040 A correspond à 6,0 V.}
\end{corrige}
\end{minipage}
\par\noindent\begin{minipage}{\linewidth}
\begin{corrige}{9}{Comparer sous la même tension}
\begin{questions}
\item $0{,}060\unite A$ puis $0{,}020\unite A$.
\item La plus faible résistance, $100\unite\Omega$, sous cette même tension.
\end{questions}
\end{corrige}
\end{minipage}
\par\noindent\begin{minipage}{\linewidth}
\begin{corrige}{10}{Doubler U}
\begin{questions}
\item $30\unite{mA}$ : tension et intensité sont proportionnelles.
\item Le dipôle reste ohmique avec la même résistance, notamment sans variation importante de température.
\end{questions}
\end{corrige}
\end{minipage}
\par\noindent\begin{minipage}{\linewidth}
\begin{corrige}{11}{Une lampe courbe}
\begin{questions}
\item Non : le quotient U/I varie avec le point de fonctionnement.
\item Lire ou mesurer le couple de valeurs au fonctionnement étudié, sans imposer une proportionnalité globale.
\end{questions}
\par\smallskip\noindent\textbf{Pour comprendre.}\enspace Pour l’exemple dessiné, U/I vaut 10 ohms au point (0,050 A ; 0,50 V), mais 25 ohms au point (0,200 A ; 5,0 V). Ce quotient variable interdit de décrire toute la courbe par une unique résistance constante.\par
\pcFigureEx{figures/cor-lampe-courbe.pdf}{La pente globale n’est pas constante ; utiliser le point de fonctionnement, sans imposer une résistance constante.}
\end{corrige}
\end{minipage}
\par\noindent\begin{minipage}{\linewidth}
\begin{corrige}{12}{Un point suspect}
\begin{questions}
\item $U=100\times0{,}030=3\unite V$.
\item La transcription, les calibres, les branchements et l’échauffement ; la donnée de 5 V est incompatible avec ce modèle constant.
\end{questions}
\par\smallskip\noindent\textbf{Pour comprendre.}\enspace Le modèle de 100 ohms prévoit $U=100\times0{,}030=3\unite V$. Le relevé 5 V est à vérifier : appareil, unité, transcription ou conditions de température peuvent être en cause.\par
\pcFigureEx{figures/cor-ohm-point-suspect.pdf}{La valeur de 5 V à 0,030 A est comparée à la prévision du modèle de 100 ohms, pas effacée sans enquête.}
\end{corrige}
\end{minipage}
\par\noindent\begin{minipage}{\linewidth}
\begin{corrige}{13}{L’ohmmètre}
\begin{questions}
\item L’appareil effectue sa propre mesure ; une tension extérieure perturbe le résultat et peut l’endommager.
\item D’autres trajets électriques peuvent fausser la résistance mesurée par rapport au composant seul.
\end{questions}
\pcFigureEx{figures/cor-ohmmetre.pdf}{Résistance déconnectée de toute alimentation ; isoler si nécessaire le dipôle des autres branches.}
\end{corrige}
\end{minipage}
\par\noindent\begin{minipage}{\linewidth}
\begin{corrige}{14}{Une tolérance fournie}
\begin{questions}
\item Oui : $209\leq225\leq231$.
\item $I=4{,}5/225=0{,}020\unite A=20\unite{mA}$.
\end{questions}
\end{corrige}
\end{minipage}
\par\noindent\begin{minipage}{\linewidth}
\begin{corrige}{15}{Un choix raisonné}
\begin{questions}
\item $R\geq U/I=6/0{,}020=300\unite\Omega$.
\item Oui : $I=6/330\approx18{,}2\unite{mA}$, inférieur à 20 mA. Il faut aussi vérifier la puissance et les tolérances.
\end{questions}
\end{corrige}
\end{minipage}
\par\noindent\begin{minipage}{\linewidth}
\phantomsection\label{cor-defi-16}\pcSolutionDefi{16}{Identifier une boîte}\begin{questions}
\item Les quotients valent $200\unite\Omega$ : comportement compatible avec un conducteur ohmique dans la plage mesurée.
\item $I=5/200=0{,}025\unite A=25\unite{mA}$.
\item La conclusion ne vaut que dans le domaine étudié ; tension, puissance, échauffement et sécurité imposent des limites. L’élève reste sur le matériel scolaire autorisé.
\end{questions}
\pcFigureEx{figures/cor-ohm-boite.pdf}{La prévision à 5 V reste dans la plage mesurée ; elle ne justifie pas une extrapolation à 100 V.}
\end{minipage}
\par\noindent\begin{minipage}{\linewidth}
\phantomsection\label{cor-activite-16}\pcSousTitre{Pistes pour l’investigation}
Une droite passant près de l’origine et des quotients U/I proches sont compatibles avec le modèle ohmique dans la plage testée. Avec I en A, R est en ohms. Un échauffement peut modifier R ; une courbe non linéaire exige un autre modèle.\par\end{minipage}
